\section{Query-Specific Window Learning}
\subsection{Two true relations, two different answers}
Let $x$ be a recording, $q$ a query, and $G$ its target interval set. RelTwin reuses two ESC-50 excerpts~\cite{piczak2015esc}, $A$ and $B$, to form two windows $W_+=(A\!\rightarrow\!B)$ and $W_-=(B\!\rightarrow\!A)$ in the same recording. Each window spans the \emph{entire ordered sequence}, not only one constituent event. For the two inverse queries, the answer-label matrix is
\begin{equation}
Y=\begin{array}{c|cc}
 &W_+&W_-\\ \hline
q_+&1&0\\
q_-&0&1
\end{array}.
\label{eq:labels}
\end{equation}
Recording-level presence labels mark both queries true but do not specify Eq.~\ref{eq:labels}. The distinction concerns supervision, not the capabilities of globally pretrained models.

Each recording has two templates, ``followed by'' and ``and then,'' in both directions. Category queries target all occurrences of $A$ or $B$ for rehearsal. Alternating the earlier relation controls a constant-position preference; excerpt reuse fixes source material within a pair. Layout-specific silence remains a potential shortcut, tested in Section~\ref{sec:timing}.

\subsection{Contrast timestamp answers during training}
For the timestamp token sequence $y(S)$ describing interval set $S$, define its mean token log-likelihood
\begin{equation}
s_\theta(x,q,S)=\frac{1}{|y(S)|}\sum_t\log P_\theta(y_t\mid x,q,y_{<t}).
\end{equation}
Candidates are $C=(W_+,W_-)$. Their normalized preferences are $p_q(j)=\operatorname{softmax}_j(s_\theta(x,q,C_j)/\tau)$, not calibrated event-presence probabilities. The paired candidate loss is
\begin{equation}
\mathcal L_{\rm cand}=-\tfrac12[\log p_{q_+}(1)+\log p_{q_-}(2)].
\label{eq:candidate}
\end{equation}
This standard candidate cross-entropy compares the correct answer with the other \emph{present} window. We add correct-answer sequence supervision and category-query rehearsal:
\begin{equation}
\mathcal L=\mathcal L_{\rm seq}+\lambda\mathcal L_{\rm cand}
+\beta\mathcal L_{\rm replay}.
\label{eq:loss}
\end{equation}
SFT uses identical data and rehearsal but omits $\mathcal L_{\rm cand}$. Equation~\ref{eq:loss}, \emph{RelTwin-Cand}, is the existing no-exchange configuration, not a newly tuned checkpoint. Inference generates intervals without access to $W_+$ or $W_-$.

\subsection{Testing whether extra objectives are needed}
Relation-Boundary Exchange Equivariance (RBEE) adds $\lambda\JS(p_{q_+},\Pi p_{q_-})$, where $\Pi$ swaps candidate indices. This enforces exchange consistency, not architectural equivariance. Perfect candidate labels already satisfy this exchange.

Setwise Preference Optimization (SetPO) adds 64 updates with six answers: both windows, their expanded versions, their union, and one spanning interval. Soft targets average set-IoU and symmetric best-match interval F1 (not one-to-one event F1). We compare it with RBEE continued for 64 updates to isolate the objective from additional training.

\subsection{Check both localization and relation selection}
For interval unions, $\setiou(S,G)=|S\cap G|/|S\cup G|$. PairAcc requires both inverse-query predictions to achieve set-IoU $\geq0.5$. We additionally require each to overlap its own target strictly more than the opposite target, yielding JointPairAcc. Swap error means that at least one query violates this strict preference, including ties.

Returning the union of two disjoint equal-length targets for \emph{both} queries gives IoU $0.5$ each: PairAcc passes but JointPairAcc fails. This task-specific diagnostic follows paired compositional evaluation~\cite{thrush2022winoground}; it complements standard IoU rather than replacing it.
